% !TeX root = ../main.tex
\section{Нормальные подгруппы, гомоморфизмы и факторгруппы}

\task{58.4, в}{Нормальные подгруппы $S_4$}
\utv{}{Все нормальные подгруппы $S_4$:
 $\{e\}$, $V_4$, $A_4$, $S_4$.
 Нетривиальные собственные --- $V_4$ и $A_4$.}{
 $A_4=\ker(\sgn)$ нормальна.
 Сопряжение перестановки переименовывает точки в её циклах;
 поэтому три двойные транспозиции переходят друг в друга,
 и $V_4\trianglelefteq S_4$.

 Докажем, что других нет. В $S_4$ классы сопряжённости
 определяются типом циклов и имеют размеры
 \[
 \begin{array}{c|ccccc}
 \text{Тип}&1^4&2\,1^2&2^2&3\,1&4\\\hline
 \text{Размер}&1&6&3&8&6
 \end{array}
 \]
 Транспозиций $\binom42=6$; разбиений четырёх точек на пары --- $3$;
 тройных циклов --- $\binom43\cdot2=8$; четверных --- $(4-1)!=6$.
 Нормальная подгруппа является объединением целых классов
 и содержит единицу. Её размер имеет вид
 $1+6a+3b+8c+6d$, $a,b,c,d\in\{0,1\}$,
 и должен делить $24$.
 Проверка этих сумм даёт только $1,4,12,24$;
 размер $4$ достигается при добавлении класса двойных транспозиций,
 размер $12$ --- при добавлении также тройных циклов.
 Это именно $V_4$ и $A_4$.
}

\task{58.5}{Нормальность не транзитивна}
Возьмём
\[G=A_4,\qquad H=V_4,\qquad
 K=\{e,(12)(34)\}.\]
Так как $V_4$ абелева, $K\trianglelefteq H$.
Из задачи 58.4 следует $H\trianglelefteq G$.
Но при $g=(123)$ имеем
\[g(12)(34)g^{-1}=(23)(14)\notin K.\]
Поэтому $K$ не нормальна в $G$.

\task{58.12}{Все нормальные подгруппы $D_n$}
Пишем
\[D_n=\langle r,s\mid r^n=s^2=e,\ srs=r^{-1}\rangle,
 \qquad |D_n|=2n.\]
Здесь $r$ --- вращение, $s$ --- отражение.
\thm{}{Нормальные подгруппы $D_n$ --- это
 \begin{itemize}
 \item все $\langle r^d\rangle$, где $d\mid n$;
 \item вся группа $D_n$;
 \item при чётном $n$ две подгруппы индекса $2$:
 $\langle r^2,s\rangle$ и $\langle r^2,rs\rangle$.
 \end{itemize}
}{
 Все подгруппы циклической группы вращений имеют вид
 $\langle r^d\rangle$. Сопряжение вращением сохраняет их,
 а сопряжение отражением заменяет $r^d$ на $r^{-d}$;
 поэтому все они нормальны.

 Теперь пусть $H$ содержит отражение $r^js$.
 Положим $H\cap\langle r\rangle=\langle r^d\rangle$.
 Произведение любых двух отражений есть вращение,
 поэтому все отражения из $H$ образуют единственный класс
 $\langle r^d\rangle r^js$ и
 $H=\langle r^d,r^js\rangle$.
 Сопряжение $r$ действует так:
 \[r(r^js)r^{-1}=r^{j+2}s.\]
 Если $H$ нормальна, отражения $r^{j+2}s$ и $r^js$
 принадлежат $H$, а их произведение равно $r^2$.
 Значит $r^2\in\langle r^d\rangle$, то есть $d\mid2$.
 При $d=1$ получаем $H=D_n$.
 При $d=2$ число $n$ чётно, а $j$ можно выбрать равным $0$ или $1$;
 получаем две указанные подгруппы порядка $n$.
 Они нормальны, поскольку имеют индекс $2$.
}
\editnote{Вложенная диэдральная подгруппа не становится нормальной
 только потому, что её порядок связан с делителем $n$.
 Например, в $D_3$ подгруппа одного отражения порядка $2$ не нормальна.}

\task{58.16}{Бесконечное ядро $\C^*\to(\R,+)$}
\utv{}{Ядро любого гомоморфизма $\varphi:\C^*\to(\R,+)$ бесконечно.}{
 Для любого корня единицы $z$ существует $m\ge1$ с $z^m=1$.
 Тогда $m\varphi(z)=\varphi(z^m)=\varphi(1)=0$.
 В аддитивной группе $\R$ нет ненулевого элемента конечного порядка,
 поэтому $\varphi(z)=0$.
 В ядре лежат все корни единицы, а их бесконечно много:
 $e^{2\pi i/m}$, $m\ge2$, попарно различны.
 Непрерывность гомоморфизма не требуется.
}

\task{58.28}{Гомоморфизмы циклических групп}
\utv{}{Все гомоморфизмы $\Z_m\to\Z_n$ имеют вид
 $\varphi_a(\overline k)=\overline{ka}$, где $ma\equiv0\pmod n$.
 Их число равно $\gcd(m,n)$.}{
 Образ $\overline1$ определяет гомоморфизм и равен некоторому $a$.
 Необходимо $0=\varphi(m\overline1)=ma$.
 Это условие также достаточно: замена $k$ на $k+mt$
 меняет $ka$ на $ka+tma$, то есть не меняет класс в $\Z_n$.
 Если $d=\gcd(m,n)$, условие $n\mid ma$ равносильно
 $n/d\mid a$. По модулю $n$ таких классов ровно $d$.
}
В разобранных пунктах получаем:
\[
\begin{array}{c|c|c}
\text{Пункт}&\text{Группы}&\text{Возможные }a=\varphi(1)\\\hline
\text{б}&\Z_6\to\Z_{18}&0,3,6,9,12,15\\
\text{г}&\Z_{12}\to\Z_{15}&0,5,10
\end{array}
\]

\task{58.29}{Гомоморфизмы $\Q\to\Z$}
\utv{}{Любой гомоморфизм $(\Q,+)\to(\Z,+)$ нулевой;
 в частности, сюръективного гомоморфизма нет.}{
 Пусть $\varphi:\Q\to\Z$, $q\in\Q$ и $t\ge1$.
 Тогда $\varphi(q)=t\varphi(q/t)$, значит целое число
 $\varphi(q)$ делится на любое положительное $t$.
 Ненулевое целое число этим свойством не обладает:
 достаточно взять $t>|\varphi(q)|$.
 Поэтому $\varphi(q)=0$ для всех $q$.
}
\editnote{Нулевой гомоморфизм существует. Условие «отобразить на»
 означает сюръективность, а не отсутствие любых гомоморфизмов.}

\task{58.30}{Факторгруппы}
\textbf{б)} $U_{12}/U_3\cong U_4\cong\Z_4$.
Положим $\zeta=e^{2\pi i/12}$ и зададим
\[\varphi:U_{12}\to U_4,\qquad\varphi(z)=z^3.\]
Это гомоморфизм. Образ $\zeta$ равен $i$ и порождает $U_4$,
поэтому отображение сюръективно.
Ядро состоит из $z$ с $z^3=1$, то есть равно $U_3$.
Теперь применяем теорему о гомоморфизме.
\editnote{Степень $4$, записанная в исходнике, даёт ядро $U_4$
 и образ $U_3$; для требуемого фактора нужна степень $3$.}

\textbf{г)} В исходнике и приложенном издании напечатано
$\R^*/\Z_+$. Если $\Z_+$ понимать как положительные целые числа
с умножением, это не подгруппа $\R^*$: например, $2^{-1}=1/2$
в неё не входит. Такая факторгруппа не определена.
Естественное исправление знаменателя --- $\Rp$. Тогда
\[\R^*/\Rp\cong\{1,-1\}\cong\Z_2.\]
Изоморфизм индуцируется гомоморфизмом знака
$\sgn:\R^*\to\{1,-1\}$ с ядром $\Rp$.
Классы --- положительные и отрицательные числа.
Это исправленная постановка, а не утверждение о буквальном
выражении $\R^*/\Z_+$.

В условии на той же странице также указаны пункты
\[\text{а)}\quad\Z/n\Z\cong\Z_n,
\qquad\text{в)}\quad4\Z/12\Z\cong\Z_3.\]
Во втором изоморфизм задаётся $4k+12\Z\mapsto k+3\Z$.

\task{58.31}{Фактор по линейному подпространству}
\utv{}{Если $H$ --- $k$-мерное подпространство $F^n$, то
 как аддитивные группы $F^n/H\cong F^{n-k}$.}{
 Выберем базис $h_1,\dots,h_k$ пространства $H$
 и дополним его до базиса
 $h_1,\dots,h_k,v_{k+1},\dots,v_n$ пространства $F^n$.
 Проекция
 \[\varphi\left(\sum_{i=1}^k a_ih_i+
                  \sum_{j=k+1}^n b_jv_j\right)
 =(b_{k+1},\dots,b_n)
 \]
 линейна, сюръективна и имеет ядро $H$.
 Применяем теорему о гомоморфизме.
 При $k=n$ пространство $F^0$ и факторгруппа тривиальны.
}
\editnote{Выражение $u|_H$ в исходнике не задаёт нужной проекции:
 вектор нельзя «ограничить на подпространство» как функцию.}

\task{58.33}{Факторы матричных групп по условиям на определитель}
Пусть
\[
\begin{gathered}
G=\GL_n(\R),\quad H=\GL_n(\C),\quad
P=\SL_n(\R),\quad Q=\SL_n(\C),\\
A=\{X\in G:|\det X|=1\},\quad
B=\{X\in H:\det X\in\R^*\},\\
C=\{X\in H:|\det X|=1\},\quad
D=\{X\in G:\det X>0\}.
\end{gathered}
\]
Все шесть факторгрупп из приведённого условия получаются
из следующей таблицы. В PDF подробно разбираются а, г, е.
\begin{center}
\begin{tabular}{c|c|c|c}
Пункт & Гомоморфизм & Ядро & Факторгруппа\\\hline
а & $G\to\R^*,\ X\mapsto\det X$ & $P$ & $G/P\cong\R^*$\\
б & $H\to\C^*,\ X\mapsto\det X$ & $Q$ & $H/Q\cong\C^*$\\
в & $G\to\{\pm1\},\ X\mapsto\sgn\det X$ & $D$ & $G/D\cong\Z_2$\\
г & $H\to U,\ X\mapsto\det X/\overline{\det X}$ & $B$ & $H/B\cong U$\\
д & $G\to\Rp,\ X\mapsto|\det X|$ & $A$ & $G/A\cong\Rp$\\
е & $H\to\Rp,\ X\mapsto|\det X|$ & $C$ & $H/C\cong\Rp$
\end{tabular}
\end{center}
Обоснуем свойства отображений. Определитель мультипликативен,
как и его модуль и знак. Для пункта г также
\[\frac{zw}{\overline{zw}}=
 \frac z{\overline z}\frac w{\overline w}.
\]
Условие $z/\overline z=1$ равносильно $z\in\R^*$,
поэтому ядро в этом пункте именно $B$.
Сюръективность определителя следует из матриц
$\diag(z,1,\dots,1)$. Для модуля выбираем $z=t>0$;
для знака --- $z=\pm1$.
Для пункта г, если $u=e^{i\theta}\in U$, выбираем
$z=e^{i\theta/2}$ и ту же диагональную матрицу:
$z/\overline z=u$.
Ядра читаются непосредственно из таблицы; отсюда также
следует нормальность всех знаменателей.
\editnote{Нормирование $X$ делением на «модуль матрицы» не задаёт
 нужного гомоморфизма. Кроме того,
 $\det(X/|\det X|)=\det X/|\det X|^n$,
 поэтому такая операция обычно не приводит определитель к модулю $1$.}
